
\(D^+(X_i)\); plaçons nous par exemple sur \(D^+(X_1)\). Alors \(p^*(\underline{E})\) est défini par le module \(\mathbb Z[T_1,\ldots,T_d,X_2/X_1,\ldots,X_d/X_1]\otimes E\), et \(\underline{O}_P(1)\) par \(\mathbb Z[T_1,\ldots,T_d,X_2/X_1,\ldots,X_d/X_1]\), le morphisme canonique envoyant \(e_i\) sur \(X_i/X_1\) pour tout \(i\). Alors \(F\) admet pour base les éléments \(e_i-e_1.X_i/X_1\) pour \(i>1\).

On a \(v(e_i-e_1.X_i/X_1)=T_i-T_1.X_i/X_1\) ; posons \(T'_i=T_i-T_1X_i/X_1\).
Comme la suite des \(T_i\) est régulière dans \(\mathbb Z[T_1,\ldots,T_d,X_2/X_1,\ldots,X_d/X_1]\), il en est de même de la suite \(T_1,T'_2,\ldots,T'_d\) qui s'en déduit par une matrice inversible, et donc aussi de la suite \(T'_2,\ldots,T'_d\) qui en est extraite (on est dans un cas noethérien). Par suite, \(v\) est un morphisme régulier.

Il faut donc montrer que les \(T'_i\) engendrent l'Idéal de \(j\). Dans \(\bigoplus_{n\ge0}J^n\), nous noterons, pour \(m\le p\), \(x^{(m)}\) l'élément \(x\) de \(J^p\), considéré comme étant de degré \(m\). Au-dessus de \(D^+(X_1)\), \(X'=\Spec\left(\left(\bigoplus_{n\ge0}J^n\right)_{(T_1^{(1)})}\right)\), et \(j\) est défini par le morphisme qui envoie \(X_i\) sur \(T_i^{(1)}\). Dans ce morphisme, l'image de \(T'_i\) est \(T_i^{(0)}-T_1^{(0)}.T_i^{(1)}/T_1^{(1)}=\bigl((T_iT_1)^{(1)}-(T_1T_i)^{(1)}\bigr)/T_1^{(1)}=0\). Donc les \(T'_i\) appartiennent à l'Idéal de \(j\).

Comme tout élément de \(\mathbb Z[T_1,\ldots,T_d,X_2/X_1,\ldots,X_d,X_1]\) peut s'exprimer comme polynôme par rapport à \(T_1,T'_2,\ldots,T'_d,X_2/X_1,\ldots,X_d/X_1\), il suffit, pour montrer que les \(T'_i\) engendrent l'Idéal de \(j\), de montrer qu'aucun élément non nul de \(\mathbb Z[T_1,X_2/X_1,\ldots,X_d/X_1]\) n'appartient à l'Idéal de \(j\), ce qui revient à montrer que l'homomorphisme
\[
\mathbb Z[T_1,X_1,\ldots,X_d]\longrightarrow\bigoplus_{n\ge0}J^n
\]
est injectif, car \(T_1^{(1)}\) n'est pas diviseur de 0 dans \(\bigoplus_{n\ge0}J^n\). Si on gradue \(\mathbb Z[T_1,X_1,\ldots,X_d]\) en attribuant le degré \(0\) à \(T_1\), et le degré \(1\) aux \(X_i\), le noyau de l'homomorphisme considéré est un idéal gradué. Or il est clair que
